\documentclass[12pt, a4paper]{article}
\usepackage[utf8]{inputenc}
\usepackage[ngerman]{babel}
\usepackage{amsmath, amssymb}
\usepackage{tikz}
\usepackage{geometry}
\geometry{margin=2.5cm}
\usepackage{parskip}

\title{\textbf{Aufgabe 2 -- Lösung} \\ Methoden I: Mathematik, 2026S}
\author{}
\date{}

\begin{document}
\maketitle

\section*{Angabe}
Seien $A$, $B$ und $C$ beliebige Mengen. Zeichnen Sie die folgenden Mengen in Venn-Diagramme ein. Was fällt Ihnen dabei auf?

\begin{itemize}
    \item $A \cap (B \cup C)$
    \item $(A \cap B) \cup (A \cap C)$
\end{itemize}

\section*{Lösung}

\subsection*{Was bedeuten die Symbole nochmal?}
\begin{itemize}
    \item $\cap$ = \textbf{Schnittmenge} ("`UND"' -- nur was in BEIDEN Mengen ist)
    \item $\cup$ = \textbf{Vereinigung} ("`ODER"' -- alles was in mindestens einer Menge ist)
\end{itemize}

\subsection*{Menge 1: $A \cap (B \cup C)$}

\textbf{Schritt für Schritt:}
\begin{enumerate}
    \item Zuerst $B \cup C$ bestimmen: Alles was in $B$ oder $C$ (oder beiden) ist.
    \item Dann davon nur das nehmen, was auch in $A$ ist (Schnitt mit $A$).
\end{enumerate}

\textit{In Worten: "`Alles was in A ist UND gleichzeitig in B oder C (oder beiden)."'}

\begin{center}
\begin{tikzpicture}[scale=0.9]
    \begin{scope}[fill opacity=0.3]
        % Highlight the intersection A ∩ (B ∪ C)
        \clip (0,0) circle (2cm);
        \fill[blue] (1.5,-1) circle (2cm);
        \fill[blue] (3,0) circle (2cm);
    \end{scope}

    % Draw circles
    \draw (0,0) circle (2cm) node[above left=1.2cm] {$A$};
    \draw (3,0) circle (2cm) node[above right=1.2cm] {$B$};
    \draw (1.5,-1) circle (2cm) node[below=1.5cm] {$C$};

    \node at (1.5, -4) {\textbf{$A \cap (B \cup C)$} -- blau markiert};
\end{tikzpicture}
\end{center}

\subsection*{Menge 2: $(A \cap B) \cup (A \cap C)$}

\textbf{Schritt für Schritt:}
\begin{enumerate}
    \item $A \cap B$: Alles was in $A$ UND $B$ ist.
    \item $A \cap C$: Alles was in $A$ UND $C$ ist.
    \item Vereinigung davon: Alles was in mindestens einem der beiden Schnitte liegt.
\end{enumerate}

\textit{In Worten: "`Alles was (in A und B) oder (in A und C) ist."'}

\begin{center}
\begin{tikzpicture}[scale=0.9]
    \begin{scope}[fill opacity=0.3]
        % A ∩ B
        \clip (0,0) circle (2cm);
        \fill[red] (3,0) circle (2cm);
    \end{scope}
    \begin{scope}[fill opacity=0.3]
        % A ∩ C
        \clip (0,0) circle (2cm);
        \fill[red] (1.5,-1) circle (2cm);
    \end{scope}

    % Draw circles
    \draw (0,0) circle (2cm) node[above left=1.2cm] {$A$};
    \draw (3,0) circle (2cm) node[above right=1.2cm] {$B$};
    \draw (1.5,-1) circle (2cm) node[below=1.5cm] {$C$};

    \node at (1.5, -4) {\textbf{$(A \cap B) \cup (A \cap C)$} -- rot markiert};
\end{tikzpicture}
\end{center}

\subsection*{Was fällt auf?}

\begin{center}
\fbox{%
    \parbox{0.8\textwidth}{%
        \centering
        \textbf{Beide Mengen sind identisch!}

        \vspace{0.5em}
        $A \cap (B \cup C) = (A \cap B) \cup (A \cap C)$

        \vspace{0.5em}
        Das ist das \textbf{Distributivgesetz} der Mengenlehre.
    }%
}
\end{center}

\textbf{Warum?} $\cap$ und $\cup$ verhalten sich wie Multiplikation und Addition:

Genau wie $a \cdot (b + c) = a \cdot b + a \cdot c$ in der Arithmetik, gilt:

$A \cap (B \cup C) = (A \cap B) \cup (A \cap C)$

\textit{Der Schnitt "`verteilt sich"' über die Vereinigung -- genau wie Multiplikation über Addition.}

\end{document}
